\chapter{Conclusions and Future Work}
\label{sec:conclusions}

This thesis has explored the application of machine learning techniques to the challenge of discovering novel polymer materials for pharmaceutical adsorption in wastewater. The work addressed both the generative task of creating new polymer structures and the predictive task of estimating their effectiveness, culminating in an integrated computational pipeline for polymer discovery.

\section{Summary of Contributions}
\label{sec:summary}

The primary contributions of this thesis are as follows.

First, a systematic data acquisition pipeline was developed to extract experimental adsorption data from the scientific literature. Despite the inherent scarcity of standardized datasets in this domain, the resulting PDCC dataset provides a foundation for machine learning experiments. Data augmentation through linear interpolation and origin point addition was shown to improve model performance, enabling more effective utilization of the limited available data. \new{As a continuation of this effort (Chapter 6), an automated, agent--based paper--scraper was implemented to mine open--access literature with several Large Language Models, yielding both a new augmentation dataset and a controlled comparison showing that the choice of model and input modality strongly affects extraction yield.}

Second, a generative model based on the minGPT architecture was successfully adapted for polymer structure generation. The model learns the grammatical rules of P-SMILES notation and produces novel, chemically valid polymer structures. Evaluation showed that 57.6\% of generated P-SMILES strings were valid, and 100\% of valid generations were novel relative to the training set. This generative capability enables the exploration of a vast chemical space that would be impractical to sample through laboratory synthesis alone.

Third, a predictive model based on the Multi-Layer Perceptron was developed to estimate adsorption capacity from chemical features. Systematic hyperparameter optimization using Leave-One-Out Cross-Validation identified configurations achieving a Q² score of approximately 0.984. While data scarcity imposes fundamental limitations, this result demonstrates that meaningful predictive patterns can be extracted from the available data. \new{As shown in Chapter 6, however, this score reflects interpolation within known polymer--molecule pairs rather than extrapolation to new ones: under stricter grouped and fixed--test evaluation the model does not generalise to unseen pairs, and augmenting the training set with the LLM--scraped data did not robustly improve this.}

Fourth, a heuristic filtering framework based on established chemical principles was implemented. Filters based on logP, TPSA, FMO theory, and synthetic accessibility effectively eliminate candidates that are unlikely to be effective or practical adsorbents. The integration of these filters with the predictive model provides a balanced approach that combines domain knowledge with data-driven prediction.

Fifth, the individual components were integrated into a complete computational pipeline that, given a target pharmaceutical molecule, generates candidate polymers, filters them based on chemical viability, and prioritizes them by predicted adsorption capacity. This pipeline represents a significant step toward computational acceleration of polymer discovery for environmental remediation.

\section{Limitations and Challenges}
\label{sec:limitations}

Several limitations of the current work must be acknowledged and addressed in future research. Data scarcity remains the most critical challenge, as the PDCC dataset contains only a few hundred experimental observations, limiting both model complexity and generalization capability. \new{Crucially, this scarcity is a property of the published literature itself, not merely of our collection effort: as the automated mining of Chapter 6 showed, even an exhaustive scrape of thousands of papers yields only a few hundred usable --- and often noisy --- measurements, because standardized polymer--molecule--capacity data are simply not yet reported at scale.} Publication bias in the training data may skew predictions toward positive outcomes, and the approximation of polymer features from capped monomers neglects higher-order structural properties that influence adsorption behavior. Additionally, the focus on homopolymers and the P-SMILES representation may restrict the applicability of the framework to more complex real-world scenarios. A more detailed discussion of these limitations is provided in Section \ref{sec:model_limitations}.

\section{Future Research Directions}
\label{sec:future_directions}

Based on the findings and limitations of this work, several promising directions for future research are identified.

The most immediate priority is data expansion. The creation of a comprehensive, publicly available database of polymer-molecule adsorption experiments would benefit not only this work but the broader research community. \new{The automated paper--scraper of Chapter 6 is a concrete step in this direction, but the re--evaluation reported there shows it is not sufficient on its own: adding hundreds of mined rows did not improve extrapolation, because the limiting factor is the number of \emph{distinct} polymer--molecule pairs rather than the number of rows per pair. The clearest path forward is therefore to grow the number of \emph{distinct, validated} polymer--molecule pairs --- well beyond the current seventy or so --- rather than the number of rows per pair, combining the automated scraper with manual revision to add quality--controlled diversity, and then to re--run the augmentation study under that larger baseline.}

\new{A second, complementary direction concerns the quality of the automatically extracted data. The cross--model comparison of Chapter 6 showed that the choice of extraction model, the input modality, the prompt, and a final manual revision all strongly affect how much trustworthy data is recovered: DeepSeek extracted far more than Gemma, image input outperformed text for the same model, and a manual pass with Claude Opus was required to resolve range--valued measurements. It is therefore reasonable to expect that stronger extraction models, better prompts, and targeted manual review will yield both more and higher--quality data. As more capable and more efficient models continue to appear, the entire scraping and augmentation study can simply be re--run, at low cost, to test whether better--quality data finally translates into improved predictive generalization.}

Model architecture exploration represents another productive direction. While this thesis focused on MLPs for prediction, Graph Neural Networks (GNNs) are naturally suited for molecular property prediction because they operate directly on molecular graphs without the need for hand-crafted feature engineering. GNNs have shown state-of-the-art performance in molecular property prediction benchmarks and could potentially improve predictive accuracy for adsorption capacity.

The generative model could be conditioned on target properties to enable directed generation. Instead of generating polymers unconditionally and then filtering, a conditional generative model could be trained to produce polymers with specific property ranges (e.g., high logP, low SA score, specific FMO gap). This approach would improve the efficiency of the discovery pipeline by focusing generation effort on chemically promising regions of the space.

The predictive model could be extended to multi-task learning, simultaneously predicting multiple adsorption metrics such as capacity, kinetics, and selectivity. Multi-task learning leverages shared representations across related tasks and can improve generalization on each individual task, particularly when data are scarce for each task individually.

Finally, experimental validation of the computationally predicted candidates is essential to close the loop between simulation and reality. The highest-ranked candidates from the discovery pipeline should be synthesized and tested under controlled laboratory conditions to verify the accuracy of the predictions and to identify systematic discrepancies that could guide model improvement.

\section{Closing Remarks}
\label{sec:closing}

This thesis has demonstrated the feasibility of applying machine learning techniques to the discovery of polymer adsorbents for pharmaceutical removal from wastewater. The integration of generative modeling, chemical filtering, and predictive analytics provides a computational framework that can accelerate the identification of promising candidates and guide experimental efforts. While data scarcity remains a fundamental constraint, the results obtained show that meaningful patterns can be learned from the available data. \new{As more standardized adsorption data are published over time, the curated dataset will grow accordingly, and the augmentation study of Chapter 6 can be revisited; the framework developed here was designed precisely to be re--run as better data and better models become available.}

The broader implication of this work is the demonstration that computational approaches can complement traditional laboratory methods in materials discovery. As datasets grow and models improve, the vision of computationally designing materials for specific environmental applications becomes increasingly achievable. The framework developed here provides a foundation upon which future advances in data, models, and experimental validation can be built, ultimately contributing to the goal of sustainable water treatment through advanced materials science.
